\subsection{Passing from the frame's finite coefficients to the arithmetic length}
\label{sec:host-selmer-length}

% BEGIN VLT1
\paragraph{VLT1. The strict Selmer group on the chosen quotient.}
\label{claim:VLT1}
Let $\Omega$ be the host's algebraic closure of $\mathbb Q$, $\ell$ its
prime, $O=c.O$ its coefficient DVR, $N>0$ its level, $e$ its exponent,
and $\psi$ its Dirichlet character. Assume the source hypotheses
$\mathrm{HostData}(\ell,e,\psi,o_1)$ and
$\mathrm{IsHostFrame}(\Omega,c,e,\psi,S,A,F)$.
Write $K=\operatorname{Frac}(O)$, $M=K/O$ with the discrete topology,
$G=\operatorname{Gal}(\Omega/\mathbb Q)$ and $\Gamma=G_S$.
The set $S$ consists of the primes dividing $\ell N$. The frame's map
$\pi:G\to\Gamma$ is the quotient by the closed normal subgroup generated
by the inertia groups above primes outside $S$. For each $v\in S$ write
$D_v,I_v\subseteq G$ for the decomposition and inertia groups of the
chosen $A_v$, and $\mathrm{loc}_v=\pi|_{D_v}$. The action on $M$ is by
$\alpha=\Psi^{-1}$ on $\Gamma$, where
$\Psi\circ\pi=\epsilon_\ell^{\,e}\psi_{\rm Gal}$.
The action is continuous since the cyclotomic character is continuous,
the Dirichlet character has open kernel, and every element of $M$ is
killed by a power of the maximal ideal. RV1 identifies $M$ with the
source's $O\otimes_{\mathbb Z_\ell}(\mathbb Q_\ell/\mathbb Z_\ell)$.

Define
\[
  \mathcal S=\bigcap_{v\in S}
   \ker\!\left(H^1_{\rm cts}(\Gamma,M(\alpha))\longrightarrow
   H^1_{\rm cts}(D_v,M(\alpha\circ\mathrm{loc}_v))\right).
\]
Then $\operatorname{length}_O\mathcal S=\mathrm{Vlen}(\Omega,c,e,\psi)$.
Indeed, the source's absolute local conditions for $\mathrm{Vlen}$ are
vanishing on $I_B$ for every valuation ring $B$ of $\Omega$, and on $D_B$
when $B(\ell N)<1$. The latter rings lie over primes of $S$.
An absolute cocycle satisfying these conditions vanishes pointwise on
every inertia group outside $S$, since the character is trivial there
and every local coboundary is zero. The action on $\ker\pi$ is trivial,
so the restriction of the cocycle to this kernel is a continuous
homomorphism to the additive group of $M$. It vanishes on the generators
and their conjugates (the conjugates are again such inertia groups),
hence on the closed subgroup they generate. The cocycle therefore
descends through $\pi$; continuity follows from the quotient topology.
The same formula $(\alpha(g)-1)m$ describes global coboundaries on both
sides after inflation, and surjectivity of $\pi$ gives injectivity on
classes. The descended class restricts to zero on every selected $D_v$.

Conversely, local coboundaries on the selected groups remain
coboundaries on their conjugates. In fact if an inflated cocycle
$z(d)=(\alpha(\pi d)-1)m$ on $D_v$, then on $gD_vg^{-1}$ its primitive
is $\alpha(\pi g)m-z(g)$ by the cocycle identity. All valuation rings
over the same rational prime are conjugate in the algebraic closure.
Thus inflation gives the required condition at every $D_B$ over $S$,
and hence at its inertia. Inertia outside $S$ maps trivially under
$\pi$. A valuation on an algebraic extension of $\mathbb Q$ that is
trivial on $\mathbb Q$ is trivial everywhere, with trivial inertia and
$B(\ell N)=1$. This exhausts the absolute conditions and gives the
claimed $O$-linear isomorphism and length identity.
% END VLT1

% BEGIN VLT2
\paragraph{VLT2. Lifting a torsion class with the frame's local conditions.}
\label{claim:VLT2}
Put $T_b=\{v_\ell\}\cup F.S_b$. At a prime $p\ne\ell$ of $S$,
membership in $F.S_b$ means that the $p$-component of the
$\ell$-power part $\kappa$ of $\psi$ is nontrivial. Put
$N_v=I_v\cap\ker(\eta\circ\pi)\cap\ker(\bar\Psi\circ\pi)$ inside $D_v$.
We first record the two local consequences needed below.

If $v\notin T_b$, the cyclotomic character is unramified there.
The character $\kappa$ is also unramified since it factors away from
the $p$-part of $N$ at this place: inertia fixes all roots of unity
of order prime to $p$. Consequently $\Psi(\pi s)$ on inertia has
order prime to $\ell$. For $s\in N_v$ its residue equals one. If a
unit $u$ of prime-to-$\ell$ order $b$ has residue one, the sum
$1+u+\cdots+u^{b-1}$ reduces to the nonzero residue of $b$, so
$(u-1)(1+u+\cdots+u^{b-1})=0$ implies $u=1$. Thus $\Psi(\pi N_v)=1$.

If $v\in T_b$, there is $x\in D_v$ for which $\alpha(\pi x)-1$
is a unit of $O$. At $v_\ell$ the element of HSR3 has
$\Psi(\pi x)^2-1$ a unit. This makes $\Psi(\pi x)-1$ and hence
$\alpha(\pi x)-1$ units. For another $v\in T_b$ write $p$ for its
prime. A nontrivial $\ell$-primary character of
$(\mathbb Z/p^{v_p(N)})^\times$ implies $\ell\mid p-1$ since
$p\ne\ell$. Thus $p$ is odd. Choose an arithmetic Frobenius lift
$\sigma\in D_v$; it acts on all prime-to-$p$ roots of unity by
raising to $p$. Hence $\bar\epsilon_\ell(\sigma)=1$ and the unit
$u$ describing its action on $\mu_N$ has dyadic coordinate $p$.

Here is the dyadic calculation at this $p$. Write $t=c_2(\psi)$,
$o=o_1\ge1$, $s=v_2(p^2-1)\ge3$ and $j=v_2(o)$.
Squaring an integer $1+2^s b$ with $b$ odd raises the exact valuation
of its difference from one by one, and an odd power preserves that
valuation. Thus $d=v_2(p^{4o}-1)=s+j+1$. The host gives
$2^d\le 2o(p^2-1)<8op^2<2^t$.
The element $p^{4o}$ generates $1+2^d\mathbb Z/2^a\mathbb Z$,
where $a=v_2(N)\ge t$, by the same squaring formula for its order.
If $\psi_2(u)^{4o}=1$, the dyadic component would factor modulo
$2^d$, contradicting its conductor $2^t$. The host requires that each
odd-prime component of the prime-to-$\ell$ part $\theta_0$ of $\psi$
is killed by $2o$. HXD1 identifies its dyadic component with $\psi_2$.
It follows that $\theta_0(u)^{4o}=\psi_2(u)^{4o}\ne1$.
Reduction kills $\kappa(u)$ and is injective on prime-to-$\ell$ roots
as shown above. Thus $\bar\psi(u)^4\ne1$. Taking $x=\sigma^2$ gives
$\bar\Psi(\pi x)^2=\bar\psi(u)^4\ne1$, which proves the asserted unit
condition at this place as well.

For $0\ne a\in O$, write $L_a=O/(a)$ with action $\alpha$, and let
$j_a:L_a\hookrightarrow M$ carry $r$ to $r/a+O$; its image is $M[a]$.
The source's $F.\mathrm{ShaO}(a)$ is the $O$-span in
$H^1_{\rm cts}(\Gamma,L_a)$ of classes having a cocycle representative
zero on $\pi N_v$ for $v\notin T_b$ and a coboundary on $D_v$ for
$v\in T_b$. Its extended $O$-length is $F.\mathrm{shaLen}(a)$.
We have
\[
    \mathcal S[a]\ \subseteq\ (j_a)_*\big(F.\mathrm{ShaO}(a)\big)
    \quad\text{inside }H^1_{\rm cts}(\Gamma,M(\alpha)).
\]
To see this, choose a cocycle $z$ representing a class in $\mathcal S[a]$.
Write $az=\delta m_0$, with $\delta m(g)=(\alpha(g)-1)m$.
Divisibility of $K/O$ gives $m_0=am_1$. The cohomologous cocycle
$z-\delta m_1$ takes values in $M[a]$ and therefore uniquely has the
form $j_a y$. The map $y$ is a continuous cocycle since $j_a$ is
injective, equivariant and both coefficient modules are discrete.
For each $v$ choose a primitive of $z|_{D_v}$ and subtract $m_1$ from
it, obtaining $b_v\in M$ with
$j_a y(\pi d)=(\alpha(\pi d)-1)b_v$ for all $d\in D_v$.
Outside $T_b$ this gives $y(\pi N_v)=0$ by the first local calculation.
Inside $T_b$ evaluate at the element $x$ of the second calculation.
Multiplication by $a$ gives $(\alpha(\pi x)-1)ab_v=0$.
Since the first factor is a unit, $ab_v=0$. Hence $b_v=j_a b_v'$
with $b_v'\in L_a$, and injectivity gives $y|_{D_v}=\delta b_v'$
through $\mathrm{loc}_v$. This places $[y]$ in the defining span and
$(j_a)_*[y]=[z]$. Inclusion into the image, followed by comparison
with its source, proves
$\operatorname{length}_O\mathcal S[a]\le F.\mathrm{shaLen}(a)$.
% END VLT2

% BEGIN VLT3
\paragraph{VLT3. The literal host length comparison.}
\label{claim:VLT3}
For every natural number $V$ satisfying
$F.\mathrm{shaLen}(a)\le V$ for all $0\ne a\in O$, one has
\[
  \mathrm{Vlen}(\Omega,c,e,\psi)\le V.
\]
Indeed $\Gamma$ is a quotient of the compact group $G$ and the image
of each continuous cocycle $\Gamma\to M$ is finite. Each value is
annihilated by a power of $\ell$ since $v_O(\ell)>0$. Taking the
largest of finitely many required exponents kills the cocycle and its
class. Thus the increasing submodules $P_k=\mathcal S[\ell^k]$ have
union $\mathcal S$. VLT2 applies at every nonzero $\ell^k$, including
$\ell^0=1$, and gives $\operatorname{length}_O P_k\le V$. Choose an
index attaining the largest length among these integers. Every later
inclusion is an equality because its finite-length quotient has
length zero. The chosen $P_k$ therefore equals $\mathcal S$. VLT1
now gives the claimed bound, including $V=0$.

In particular $\mathrm{Vlen}$ is finite before the operation
$\mathrm{toNat}$ in the source's real normalization.
With $e_O=\operatorname{length}_O(O/\ell O)=v_O(\ell)>0$ from RV1,
the contribution to $\mathrm{VLe}$ consequently satisfies
\[
   \mathrm{Vn}(\Omega,c,e,\psi)\log\ell
      =\frac{\operatorname{length}_O\mathcal S}{e_O}\log\ell
      \le \frac{V}{e_O}\log\ell .
\]
Here $V$ is precisely the uniform upper bound on the host's finite
coefficient lengths. It need not be attained by one of those lengths.
% END VLT3
